\section{Market-regulation scenarios}
\label{sec:regulation}

The regulatory analysis compares policies under the same technology, initial
conditions, and population path. The central issue is that the current-deployment
distortion and the research-incentive distortion operate through different margins.

\subsection{Laissez-faire monopoly}

The integrated developer chooses the quantity and price of AI services together with
human and computational research inputs. It receives all operating rents and no
additional research support. This scenario provides the benchmark paths for AI
deployment, capability, research, monopoly rents, per-capita consumption, and the
automated share of research.

\subsection{Regulated access price}

A price cap or access rule moves the AI-service price toward marginal cost. It expands
current use of AI at a given capability but reduces operating rents and therefore the
private shadow value of frontier capability. The experiment measures the static gain
from deployment against the induced change in research and its input composition.

\subsection{Research subsidy or innovation prize}

If an explicit planner problem implies $Q_t>q_t$, a subsidy to human research or
research compute, or an innovation prize contingent on capability improvements, can
target the resulting wedge $(Q_t-q_t)F_{M,t}$. The baseline model does not establish
the sign of this wedge. The distinction matters for implementation: an input subsidy
rewards expenditure, while a prize rewards measured output. Either instrument can
increase private research without requiring a high markup on all current uses of AI.

\subsection{Joint price and research policy}

A joint policy combines regulated access to current AI services with explicit
remuneration for frontier innovation. In principle, the access rule targets the
static monopoly distortion. A research instrument would address a separate dynamic
wedge only if a fully specified planner problem establishes one. A quantitative
welfare extension would determine whether one instrument is enough or whether this
two-instrument policy is required.

\subsection{Policy comparison}

For each regime, the model will report:
\begin{enumerate}[label=(\roman*)]
    \item the path of the automated share of effective research;
    \item the paths of AI capability, aggregate output, and output per capita;
    \item consumption and factor-income shares;
    \item AI-service prices, quantities, markups, and developer rents;
    \item the level and composition of research; and
    \item aggregate welfare and welfare by household type.
\end{enumerate}
The comparison will distinguish policies that change current deployment from those
that change the private appropriation of innovation benefits. It will also show how
population growth affects the human-dominated benchmark against which an
automation-driven change in the growth regime is measured.
